\subsection{Four-parameter synthetic recovery on PAT16 anatomy}\label{sec:pat16_recovery_methods}
A third inference experiment fitted synthetic longitudinal signals on the same real preoperative PAT16 anatomy and dMRI directions. It jointly estimated separate $D_{\rm WM}$ and $D_{\rm GM}$, shared $\rho$ and one migration ratio $r=D_\parallel/D_\perp$ shared across eligible white matter in Equation~\eqref{eq:pat16_tensor}. Known generating values were $(0.500,0.100,0.040,2.500)$, with diffusion in mm$^2$/day, proliferation in day$^{-1}$ and dimensionless $r$. Mean diffusivity and anisotropy were separate parameters; no fixed tissue ratio, ordering between WM and GM diffusivities or probability prior was imposed. Gray matter and non-eligible white matter remained isotropic. Tissue geometry, directions, eligibility and normalized initialization were fixed. Uniform $K$ cancels from these normalized dynamics; neither absolute capacity nor cell counts were identified. Unlike the pilot's assumed-parameter sensitivity bands, the intervals here were fitted conditional confidence intervals. Neither experiment estimated rates from measured patient tumor signals.

The full 3 mm tissue domain was evolved, with double-precision snapshots for likelihood differentiation. The known initial density used lesion threshold 0.5, smoothing with standard deviation 0.7 grid voxels and amplitude 0.6. A prespecified 72 mm observation cube centered on the tissue-restricted initial lesion centroid supplied $12^3=1{,}728$ voxels per scan. Gaussian blur with 2 mm standard deviation preceded 6 mm voxel averaging with $3^3$ midpoint quadrature samples and physical-coordinate interpolation. Independent Gaussian noise of known standard deviation 0.015 was added after imaging, without clipping. The synthetic cellularity scale was assumed known; these signals were not T1, FLAIR or measured patient cellularity. BOLD and functional connectivity did not enter the likelihood.

The generator and inverse shared spatial anatomy, the tensor stencil, initialization and observation model, with maximum steps of 0.25 and 0.5 day, respectively. All 5,184 scalar measurements at days 20/40/60 were fitted; days 75/90 were withheld from optimization and uncertainty estimation. Three starts minimized standardized Gaussian residuals in four log parameters. Numerical bounds were $(0.05,0.02,0.01,1.0)$ to $(1.0,0.6,0.09,3.5)$; the tensor family assumes prolate migration ($r\geq1$), and the upper bound is a monotonicity restriction of this stencil. Full-rank Jacobian information supplied local marginal 95\% log-Wald intervals, using the known noise variance without residual-variance rescaling.

We compared local intervals with percentile intervals from \PRBootstraps{} parametric bootstrap refits of noisy early scans generated at the fitted parameters. A separate \PRTrials{}-dataset audit added fresh noise at the known, finer-time generating parameters to evaluate bias, RMSE and marginal local-interval coverage with exact binomial uncertainty. Every attempted trial remained in the coverage denominator; failed or boundary local intervals counted as misses. Marginal bootstrap parameter-interval coverage and joint four-parameter coverage were not audited. Bootstrap trajectories produced pointwise signal intervals. For new-scan prediction, we compared the original 100 noise-augmented trajectories with \PPDraws{} predictive draws resampling these trajectories, and with numerical quantiles of their exact Gaussian-noise mixture. The latter supplies the final prediction intervals, removing noise-tail sampling error while retaining finite-bootstrap error. All interval levels are nominal 95\%. Forecasts were evaluated over the full cube and a prespecified 24 mm sphere around its center, containing 280 observation-voxel centers. A separate \PPTrials{}-dataset predictive audit regenerated early and withheld observations independently at the known finer-time signal. Each dataset was refitted from the default start, followed by 100 new bootstrap refits and all three predictive constructions. Thus the audit refitted uncertainty in each repetition rather than testing repeated noise against one fixed band. We report coverage, width and the proper 95\% interval score \citep{gneiting2007}. Monte Carlo intervals for mean spatial coverage use independent datasets as replicates; a fixed central voxel additionally receives exact binomial coverage limits. Withheld observations were used only for evaluation; no tail or width correction was tuned against them. Noise-free matched-model recovery, a finer-time noise-free fit and a time-zero-only negative control checked implementation and local identifiability. Appendix~\ref{app:pat16_recovery} gives the operator, optimization and reproducibility details.
